\section{Loanword Adaptation as a Transducer}

We can think of the computational task of LA as a string transduction
problem

\[
    L_2 \mapsto F^*
\]

where \( F \) is the alphabet of \( \mathcal{F} \) which is the same as
the latin alphabet but with the letter \~N and Ng added. Note that ng, 
as the velar nasal \textipa{N}, is counted as a single character.

By framing it as a transduction problem, that is, given an word \( w \in L_2 \)
the transducer \( A: L_2 \mapsto F^* \) returns \( A(w) \in F^* \)
where \( A(w) \) is the adapted loanword in \( \mathcal{F} \). We can
employ formalisms used in formal language theory to perform the task.

Taking inspiration from \underline{O-COCOSDA paper} and~\cite{Mao_Hulden_2016},
we can use context-sensitive rewrite rules defined as an attribute grammar 
(AG) to implement the LA algorithm.

Since LA requires constraints from \( F \) and the
orthography and phonetics of the \( w \), we can embed these
constraints as attributes required to be satisfied to perform the rewrite.

Hence, we will have a context-sensitive grammar (CSG) \( G \), which defines
the rewrite-rules as productions of the form

\[
    \texttt{rule\_i}(\omega_i) \to 
    \omega_{j} 
    \qquad 
    \mid \omega_i \vDash C
\]

Informally, we will rewrite \( \omega_i \) to \( \omega_j \) if and only if
\( \omega_i \) satisfies each constraint in \( C \).